
\chapter{Rappels d'algèbre linéaire}

\section{Opérations sur les matrices}

\subsection{Définitions}

\begin{defn}
Une \textbf{matrice} réelle $m\times n$ est un tableau ordonné de nombres réels arrangés en $m$ lignes et $n$ colonnes. On appelle ce tableau matrice de dimension $m \times n$. \\
Notation
$$
A=\left[ \begin{array}{ccccc}
a_{11} & \cdots & a_{1j} & \cdots & a_{1n} \\
\vdots &         & \vdots &           & \vdots \\
a_{i1} &  \cdots       &  a_{ij} & \cdots & a_{in} \\
\vdots &         & \vdots &           & \vdots \\
a_{m1} &  \cdots & a_{mj}& \cdots & a_{mn}
\end{array}
\right]
$$
\end{defn}

\begin{defn}
La \textbf{somme matricielle} $C=A+B$ de deux matrices $A$ ($m\times n$) et $B$ ($m\times n$) est la matrice $C$ (de même dimension que $A$ et $B$) dont les éléments sont
$$c_{ij}=a_{ij}+b_{ij}\qquad\;\;i=1,\ldots,m; j=1,\ldots,n.$$
\end{defn}

\begin{defn}
Soit $A$ une matrice de dimension $m\times n$ quelconque, alors la \textbf{matrice transposée} de $A$, notée $A'$ de dimension $n\times m$ comprend les éléments $\alpha_{ji}=a_{ij},\;\;i=1,\ldots,m; j=1,\ldots,n.$
\end{defn}

\begin{defn}
Le produit du scalaire $k, k\in\R$, par la matrice $A$ de dimension $m\times n$ est la matrice $C$ de dimension $m\times n$, $C=kA$, dont les éléments sont
$$c_{ij}=ka_{ij}\;\;i=1,\ldots,m; j=1,\ldots,n.$$
\end{defn}

\begin{defn}
Soient $A$ et $B$ deux matrices de dimensions respectives $m\times r$ et $r\times n$. Le \textbf{produit matriciel} $C=AB$ correspond à la matrice $C$ de dimension $m\times n$ dont les éléments s'obtiennent comme
$$c_{ij}=\sum_{k=1}^r a_{ik}b_{kj}\;\;i=1,\ldots,m; j=1,\ldots,n.$$
Il est à noter que le produit matriciel n'est défini que si le nombre de colonnes de $A$ est égal au nombre de lignes de $B$. L'élément $c_{ij}$ correspond au produit scalaire euclidien du i$^\text{ème}$ vecteur ligne de $A$ par le j$^\text{ème}$ vecteur colonne de $B$.
\end{defn}



\subsection{Propriétés du calcul matriciel}
\label{sec:prop_calc_mat}

En supposant des matrices compatibles pour les opérations en question, on a que

\begin{enumerate}
\item Commutativité de la somme : $A + B  =  B + A$
\item Associativité de la somme : $(A+B)+C  =  A+ (B+C)$
\item Associativité du produit:  $(AB)C=A(BC)$
\item Divers types de distributivité :\\
$A(B+C)=AB+AC$ \\
$(A+B)C=AC+BC$ \\
$k(A+B)=kA+kB$ \\
$(k_1+k_2)A=k_1A+k_2A$
\item Divers types d'associativité : \\
$k_1(k_2C)=k_1k_2C$  \\
$k(AB)=(kA)B=A(kB)$
\end{enumerate}


\subsection{Propriétés de la transposition}
\label{sec:prop_transp}
\begin{enumerate}
\item
$(A+B)'=A'+B'$
\item
$(kA)'=kA'$
\item
$(AB)'=B'A'$
\item
$(A')'=A$
\end{enumerate}

\section{Espérance et variance de combinaisons linéaires de variables aléatoires}

 Soit $\boY=(Y_1\ Y_2\ \cdots\ Y_n)'$, un vecteur de
$n$ variables al\'eatoires et $\bo{A}$ et $\bo{B}$, des matrices de constantes. L'espérance et la variance de $\boY$ s'écrivent sous forme matricielle comme
\begin{eqnarray*}
\esp[\boY]&\equiv&\left(\begin{array}{c} \esp[Y_1]\\ \esp[Y_2]\\ \vdots \\ \esp[Y_n]\end{array}\right)\
\ \var[\boY]\equiv\left(\begin{array}{cccc}
\var[Y_1] & \cov(Y_1,Y_2) & \cdots & \cov(Y_1,Y_n) \\
\cov(Y_2,Y_1) & \var[Y_2] & \cdots & \cov(Y_2,Y_n) \\
\vdots & \vdots & \ddots & \vdots \\
\cov(Y_n,Y_1) & \cov(Y_n,Y_2) & \cdots & \var[Y_n]\end{array}\right)\\
\end{eqnarray*}
Les propriétés de l'espérance et de la variance sous forme matricielle sont
\begin{eqnarray*}
\esp[\bo{A}\boY+\bo{B}]&=&\bo{A}\esp[\boY]+\bo{B} \\
\var[\bo{A}\boY] &= & \bo{A}\var[\boY]\bo{A}' \\
\var[\bo{A}\boY+\bo{B}] & = & \bo{A}\var[\boY]\bo{A}'.
\end{eqnarray*} 